\documentclass[12pt,a4paper]{article}

\usepackage{fontspec}
\defaultfontfeatures{Ligatures=TeX}
\setmainfont[
  Path={fonts/TH Sarabun PSK V-1/},
  Extension=.ttf,
  UprightFont={*},
  BoldFont={* Bold},
  ItalicFont={* Italic},
  BoldItalicFont={* BoldItalic}
]{THSarabun}
\setmonofont{Latin Modern Mono}
\usepackage{amsmath,amssymb,mathtools}
\usepackage{geometry,booktabs,array,tabularx,longtable}
\usepackage{siunitx}
\usepackage{xcolor,hyperref,fancyhdr,caption,subcaption,float}
\usepackage{pgfplots,pgfplotstable}
\usepgfplotslibrary{groupplots}
\pgfplotsset{compat=1.18}
\geometry{margin=1in}
\setlength{\headheight}{14.5pt}
\pagestyle{fancy}\fancyhf{}
\lhead{IEEE 14-Bus SG + Four GFL-RMS10 System}
\rhead{\thepage}
\hypersetup{colorlinks=true,linkcolor=blue!55!black,urlcolor=blue!55!black,citecolor=blue!55!black}
\sisetup{detect-all,round-mode=places,round-precision=5,table-number-alignment=center}
\renewcommand{\arraystretch}{1.12}
\captionsetup[table]{skip=7pt}
\emergencystretch=2em
\linespread{1.15}

\definecolor{sgblue}{RGB}{0,114,189}
\definecolor{ibrorange}{RGB}{217,83,25}
\definecolor{ibryellow}{RGB}{237,177,32}
\definecolor{ibrpurple}{RGB}{126,47,142}
\definecolor{ibrgreen}{RGB}{119,172,48}
\pgfplotsset{
 report axis/.style={
   width=0.43\textwidth,height=0.25\textwidth,
   grid=major,grid style={gray!18},axis line style={black!65},
   tick label style={font=\small},label style={font=\small},
   title style={font=\small\bfseries},legend style={font=\scriptsize,draw=black!35,
   fill=white,legend cell align=left},line width=0.85pt,
   scaled ticks=false,unbounded coords=jump
 },
 report wide/.style={report axis,width=0.94\textwidth,height=0.34\textwidth},
 every axis plot/.append style={line join=round}
}
\newcommand{\datadir}{figures/ieee14_ibr_pf_sssa_ts}
\newcommand{\code}[1]{\texttt{\detokenize{#1}}}
\newcommand{\clamp}{\operatorname{clamp}}
\newcommand{\hold}{\operatorname{conditional\_hold}}
\newcommand{\vect}[1]{\boldsymbol{#1}}

\title{\textbf{Dynamic-Equation and Analysis Contract for an IEEE 14-Bus Mixed-Resource System}\\
\large One EMF6 Synchronous Generator and Four 10-State GFL--RMS10 Converters\\[1mm]
\normalsize Power Flow, Equilibrium, Full-State SSSA, and Event-Free Time-Domain Simulation}
\author{In-house Power-Flow Project}
\date{19 July 2026}

\begin{document}
\maketitle

\begin{abstract}
This report documents the exact positive-sequence RMS equations used by the
project for the IEEE 14-bus mixed-resource case.  The fixed operating mode
contains SG1 at bus~1 and IBR2, IBR3, IBR6, and IBR8 in 10-state
PLL-resolved grid-following (GFL) mode.  The report
defines every state, reference frame, sign convention, per-unit conversion,
network residual, equilibrium condition, SSSA reduction, and implicit
trapezoidal time step.  The reported operating-point tables, SSSA, and the
15-s event-free TS all use the same accepted SG1 plus four-GFL equilibrium.
Numerical plots are rendered natively by
\LaTeX{} from CSV files exported from the project-owned MATLAB solvers, so the
plot font, size, and spacing are identical to the report.  Numerical
instability or fail-closed outcomes are reported rather than tuned away.
\end{abstract}

\newpage
\tableofcontents
\newpage

\section{Scope, evidence discipline, and reproducibility}

The production calculation uses only project-owned base-MATLAB code.  The
allowed numerical primitives include \verb|\| and \verb|eig|; no external
power-flow, optimization, Simulink, or loaded-solution result feeds the
calculation.  The evidence files in \code{\datadir} are generated by
\code{scripts/reporting/generate_ieee14_ibr_pf_sssa_ts_report_data.m}.  That
script calls the same production dispatch used by \code{solve_case}, writes
the full-precision results to CSV and MAT files, and records branch, commit,
profile, convergence, and physical classifications in \code{manifest.txt}.

The products in this report are deliberately separated:
\begin{enumerate}
  \item mode-aware all-GFL PF warm-start with effective PQ converter buses;
  \item mixed-device equilibrium and full-state SSSA at the fixed four-GFL mode map;
  \item a 15-s event-free normal-operation TS from that equilibrium.
\end{enumerate}

\section{Case data and index maps}

\subsection{System and resource map}

The network is the 100-MVA, 60-Hz IEEE 14-bus data set.  The external bus ID
is not assumed equal to an arbitrary internal index, even though this case
happens to use consecutive IDs.  Algebraic voltage order is
\[
\vect y=[V_{x,1},V_{y,1},V_{x,2},V_{y,2},\ldots,V_{x,14},V_{y,14}]^{\mathsf T},
\qquad V_k=V_{x,k}+jV_{y,k}.
\]
The fixed normal resource map is
\begin{center}
\begin{tabular}{ccllcc}
\toprule
Device index & Bus ID & Device & Normal mode & Physical states & Active in normal SSSA\\
\midrule
1 & 1 & SG1 & synchronous EMF6 & 6 & 5\\
2 & 2 & IBR2 & GFL-RMS10 & 10 of dual 23 & 10\\
3 & 3 & IBR3 & GFL-RMS10 & 10 of dual 23 & 10\\
4 & 6 & IBR6 & GFL-RMS10 & 10 of dual 23 & 10\\
5 & 8 & IBR8 & GFL-RMS10 & 10 of dual 23 & 10\\
\midrule
\multicolumn{4}{r}{Active dynamic order} & \multicolumn{2}{c}{\(5+4(10)=45\)}\\
\bottomrule
\end{tabular}
\end{center}
The runtime inventory remains fixed at 98 states: SG1 contributes 6 and each
dual-capable IBR reserves 23 coordinates.  For this report, only the ten
GFL-RMS10 coordinates of each IBR enter the 45-state reduced dynamic matrix.

\subsection{Bus and line input data}

\begin{table}[H]\centering\scriptsize
\caption{Case bus data.  These are inputs, not solved PF results.}
\pgfplotstabletypeset[col sep=comma,string type,
 columns={bus_id,type,Pd_MW,Qd_MVAr,Vm0_pu,Va0_deg,base_kV,Vmax_pu,Vmin_pu},
 columns/bus_id/.style={column name=Bus},columns/type/.style={column name=Type},
 columns/Pd_MW/.style={column name=$P_L$ MW},columns/Qd_MVAr/.style={column name=$Q_L$ MVAr},
 columns/Vm0_pu/.style={column name=$V_0$ pu},columns/Va0_deg/.style={column name=$\theta_0$ deg},
 columns/base_kV/.style={column name=Base kV},columns/Vmax_pu/.style={column name=$V_{max}$},
 columns/Vmin_pu/.style={column name=$V_{min}$},
 every head row/.style={before row=\toprule,after row=\midrule},every last row/.style={after row=\bottomrule}]
{\datadir/case_bus_data.csv}
\end{table}

\begin{table}[H]\centering\small
\caption{Case branch series and charging data in project index order.}
\pgfplotstabletypeset[col sep=comma,
 columns={branch_index,from_bus,to_bus,r_pu,x_pu,b_pu},
 columns/branch_index/.style={column name=Index},columns/from_bus/.style={column name=From},
 columns/to_bus/.style={column name=To},columns/r_pu/.style={column name=$r$ pu},
 columns/x_pu/.style={column name=$x$ pu},columns/b_pu/.style={column name=$b_c$ pu},
 columns/r_pu/.append style={fixed,precision=5},columns/x_pu/.append style={fixed,precision=5},
 columns/b_pu/.append style={fixed,precision=4},
 every head row/.style={before row=\toprule,after row=\midrule},every last row/.style={after row=\bottomrule}]
{\datadir/case_line_data.csv}
\end{table}

\begin{table}[H]\centering\small
\caption{Case branch ratings, transformer ratios, phase shifts, and status in the same index order.}
\pgfplotstabletypeset[col sep=comma,
 columns={branch_index,from_bus,to_bus,rateA_MVA,tap,shift_deg,status},
 columns/branch_index/.style={column name=Index},columns/from_bus/.style={column name=From},
 columns/to_bus/.style={column name=To},
 columns/rateA_MVA/.style={column name=Rate A},columns/tap/.style={column name=Tap},
 columns/shift_deg/.style={column name=Shift deg},columns/status/.style={column name=Status},
 columns/rateA_MVA/.append style={fixed,precision=1},columns/tap/.append style={fixed,precision=3},
 columns/shift_deg/.append style={fixed,precision=1},
 every head row/.style={before row=\toprule,after row=\midrule},every last row/.style={after row=\bottomrule}]
{\datadir/case_line_data.csv}
\end{table}

The source bus-data base-kV column is zero (unspecified).  The project case
contract supplies a 69-kV reporting base, so solved values use
\(V_{\mathrm{kV}}=69|V|_{\mathrm{pu}}\); this reporting conversion does not
enter the per-unit PF, equilibrium, SSSA, or TS equations.

For a branch \(i\!\to\!j\), \(z_{ij}=r_{ij}+jx_{ij}\),
\(y_{ij}=1/z_{ij}\), and the charging susceptance is divided equally between
the terminal shunts.  Tap magnitude \(\tau\) and phase shift \(\phi\) enter
the off-nominal transformer terms.  These contributions assemble the complex
network matrix \(Y\).

\section{Common frames, signs, and per-unit conversion}

All devices use generator sign: positive complex power is injected into the
network,
\[
S=P+jQ=V I^*,\qquad I_{bus}=\sum_{d\in bus} I_d.
\]
For an IBR whose local angle is \(\delta_c\), the network-to-device transform is
\begin{align}
v_d&=V_x\cos\delta_c+V_y\sin\delta_c,&
v_q&=-V_x\sin\delta_c+V_y\cos\delta_c,\\
i_x&=i_d\cos\delta_c-i_q\sin\delta_c,&
i_y&=i_d\sin\delta_c+i_q\cos\delta_c.
\end{align}
Thus \(v_q=0\) when a PLL is locked to the positive-sequence bus voltage.
With this convention,
\[
P_{inv}=v_di_d+v_qi_q,\qquad Q_{inv}=v_qi_d-v_di_q.
\]
If \(S_{base}\) is the system base and \(M_{base}\) is the converter base,
\[
\kappa=\frac{S_{base}}{M_{base}},\qquad
P_{inv}=\kappa P_{sys},\quad Q_{inv}=\kappa Q_{sys},\quad
I_{sys}=\frac{I_{inv}}{\kappa}.
\]
IBR2 uses \(M_{base}=140\) MVA, hence \(\kappa=100/140\); IBR3, IBR6,
and IBR8 use 100 MVA and \(\kappa=1\).

\section{Composite network DAE}

Every device supplies a differential right-hand side and a positive complex
current injection.  The composite equations are
\begin{align}
\dot{\vect x}&=\vect f(\vect x,\vect y,\vect u;\sigma),\\
\vect 0&=\vect g(\vect x,\vect y,\vect u;\sigma)
          =\mathcal R\!\left(Y(\sigma)V-I_{bus}(\vect x,\vect y,\vect u;\sigma)\right),
\end{align}
where \(\mathcal R(z)=[\Re z_1,\Im z_1,\ldots]^{\mathsf T}\) and \(\sigma\)
contains topology, online status, and device mode.  At a declared reference
bus, the appropriate KCL rows are replaced by explicit voltage constraints;
the replacement is metadata-driven, not inferred inside a device.

\section{Synchronous generator: EMF6 equations}

\subsection{State order and meaning}

\[
\vect x_{SG}=[\delta,\omega,E'_q,E'_d,E''_q,E''_d]^{\mathsf T},\qquad
\vect u_{SG}=[T_m,E_{fd}]^{\mathsf T}.
\]
\begin{center}\small
\begin{tabularx}{\textwidth}{c l l X}
\toprule
Local & State & Unit & Physical meaning\\\midrule
1 & \(\delta\) & rad & electrical rotor angle in the network reference frame\\
2 & \(\omega\) & pu & rotor-speed deviation; actual speed is \(\omega_b(1+\omega)\)\\
3 & \(E'_q\) & pu & q-axis transient internal EMF\\
4 & \(E'_d\) & pu & d-axis transient EMF; frozen at zero here because \(T'_{q0}=0\)\\
5 & \(E''_q\) & pu & q-axis subtransient internal EMF\\
6 & \(E''_d\) & pu & d-axis subtransient internal EMF\\
\bottomrule
\end{tabularx}
\end{center}
The physical inventory is sixth order, but the normal SSSA contribution is
five because the sourced Kodsi round-rotor data impose the singular limit
\(T'_{q0}=0\Rightarrow E'_d=0\).

\subsection{Stator algebraic equations and ODEs}

The project Kundur transform gives \((V_d,V_q)\) from \((V_x,V_y,\delta)\).
Defining \(r_d=V_d-E''_d\), \(r_q=V_q-E''_q\), and
\(\Delta=R_a^2+X''_dX''_q\), the stator current is
\begin{align}
I_d&=\frac{-R_a r_d-X''_q r_q}{\Delta},&
I_q&=\frac{ X''_d r_d-R_a r_q}{\Delta},\\
T_e&=V_dI_d+V_qI_q+R_a(I_d^2+I_q^2).
\end{align}
On the system base, the connected-machine dynamics are
\begin{align}
\dot\delta&=\omega_b\omega,\\
2H_{sys}\dot\omega&=T_m-T_e-D_{sys}\omega,\\
T'_{d0}\dot E'_q&=E_{fd}+c_dE''_q-d_dE'_q,\\
T'_{q0}\dot E'_d&=c_qE''_d-d_qE'_d,\\
T''_{d0}\dot E''_q&=E'_q-E''_q-(X'_d-X''_d)I_d,\\
T''_{q0}\dot E''_d&=E'_d-E''_d+(X'_q-X''_q)I_q.
\end{align}
When the SG breaker is open, \(I_d=I_q=T_e=0\).  The angle and speed coast,
and the four flux equations use the same formulas with their current terms
removed.  The retained \(T_m\) policy is therefore visible in the post-trip
trajectory; it is not silently re-solved at every time step.

\section{GFM--VSG REGFM\_B1: 13-state equations}

These equations document the reserved grid-forming branch of each dual-mode
IBR.  That branch is inactive throughout the operating point studied in this
report: it contributes no row to the 45-state SSSA matrix and is not integrated
by the event-free TS.  It is retained here to make the complete device model
and inactive-state inventory auditable.

\subsection{State order and device frame}

\begin{multline}
\vect x_{GFM}=[\omega_m,\delta_{IT},x_w,x_E,\delta_{PLL},\xi_{PLL},
P_f,I_{d,f},Q_f,V_f,I_{q,f},\delta_{IT}^{max},\delta_{IT}^{min}]^{\mathsf T}.
\end{multline}
\begin{center}\scriptsize
\begin{tabularx}{\textwidth}{c l X}
\toprule
Local & State & Meaning\\\midrule
1 & \(\omega_m\) & VSG speed deviation\\
2 & \(\delta_{IT}\) & PLL-relative VSG internal angle\\
3 & \(x_w\) & damping washout state\\
4 & \(x_E\) & voltage-controller integral state\\
5--6 & \(\delta_{PLL},\xi_{PLL}\) & PLL angle and PI integral\\
7--11 & \(P_f,I_{d,f},Q_f,V_f,I_{q,f}\) & measured active power, d current, reactive power, voltage and q current filters\\
12--13 & \(\delta_{IT}^{max},\delta_{IT}^{min}\) & dynamic upper/lower current-angle bounds\\
\bottomrule
\end{tabularx}
\end{center}
The device voltage angle is
\[
\delta_{VSM}=\delta_{PLL}+\clamp(\delta_{IT},\delta_{IT}^{min},\delta_{IT}^{max}),
\quad
\clamp(z,l,h)=\min(\max(z,l),h).
\]
The clamp is an algebraic projection: an interior value is unchanged, a value
below \(l\) is replaced by \(l\), and a value above \(h\) is replaced by \(h\).
The derivative uses an outward-hold/inward-release rule rather than
differentiating through an impossible outward motion.

\subsection{PLL, VSG, voltage and filter dynamics}

With \(V_q=-V_x\sin\delta_{PLL}+V_y\cos\delta_{PLL}\),
\begin{align}
\dot\xi_{PLL}&=V_q,&
\dot\delta_{PLL}&=\omega_b(k_{p,PLL}V_q+k_{i,PLL}\xi_{PLL}),\\
2H\dot\omega_m&=P_{ref,inv}-P_f-(1/m_p+D_1)\omega_m-D_2(\omega_m-x_w),\\
\dot x_w&=\omega_D(\omega_m-x_w),\\
\dot\delta_{IT}&=\hold\!\left(\omega_b\omega_m-\dot\delta_{PLL};
\delta_{IT}^{min},\delta_{IT}^{max}\right),\\
\dot x_E&=\hold_E(V_{ref}-V_f),\\
E_{VSM}^{raw}&=V_{ref}-m_qQ_f+k_{pv}(V_{ref}-V_f)+k_{iv}x_E.
\end{align}
The voltage-integrator hold is active only when the voltage command is at a
limit and the error would drive it farther outward.  The measurement states are
\begin{align}
T_{pf}\dot P_f&=\kappa P-P_f,&T_{If}\dot I_{d,f}&=\kappa I_d-I_{d,f},\\
T_{Qf}\dot Q_f&=\kappa Q-Q_f,&T_{Vf}\dot V_f&=|V|-V_f,\\
T_{If}\dot I_{q,f}&=\kappa I_q-I_{q,f}.
\end{align}
The moving angle bounds satisfy
\begin{align}
\dot\delta_{IT}^{max}&=\hold\bigl(k_I(I_{d,max}^{SS}-I_{d,f});0,\delta_{max}\bigr),\\
\dot\delta_{IT}^{min}&=\hold\bigl(k_I(-K_eI_{d,max}^{SS}-I_{d,f});-\delta_{max},0\bigr),
\end{align}
with the lower state frozen when its source flag is disabled.

The unsaturated network current is the voltage-behind-impedance relation
\[
I_{unc}=\frac{E_{VSM}e^{j\delta_{VSM}}-V}{\kappa(R_e+jX_L)}.
\]
Its magnitude is circularly limited at \(I_{maxF}/\kappa\), after the sourced
P/Q and voltage-limit logic.  This current, not a separate SSSA surrogate, is
injected into the common KCL equations.

\section{GFL--RMS10: 10-state PLL-resolved equations}

\subsection{State order and meaning}

\[
\vect x_{GFL}=[\delta_{PLL},\xi_{PLL},P_f,Q_f,\xi_P,\xi_Q,
\xi_{id},\xi_{iq},i_d,i_q]^{\mathsf T}.
\]
\begin{center}\small
\begin{tabularx}{\textwidth}{c l l X}
\toprule
Local & State & Unit & Meaning\\\midrule
1 & \(\delta_{PLL}\) & rad & SRF-PLL angle used by every dq transform\\
2 & \(\xi_{PLL}\) & pu s & PLL PI integral of \(v_q\)\\
3--4 & \(P_f,Q_f\) & pu inv. & filtered measured injected powers\\
5--6 & \(\xi_P,\xi_Q\) & pu s & outer active/reactive-power PI integrals\\
7--8 & \(\xi_{id},\xi_{iq}\) & pu s & inner d/q current-controller PI integrals\\
9--10 & \(i_d,i_q\) & pu inv. & physical L-filter current states in the PLL frame\\
\bottomrule
\end{tabularx}
\end{center}
The last two rows are why this model is a genuine dynamic GFL rather than a
PLL label attached to algebraic current commands.

\subsection{PLL, power measurement, and exact P/Q feedforward}

Let \(D_V=v_d^2+v_q^2\).  The declared positive-sequence domain requires
\(D_V\ge V_{div,min}^2\).  Within that domain,
\begin{align}
\dot\xi_{PLL}&=v_q,&
\Delta\omega_{PLL}&=k_{p,PLL}v_q+k_{i,PLL}\xi_{PLL},&
\dot\delta_{PLL}&=\omega_b\Delta\omega_{PLL},\\
T_P\dot P_f&=P_{inv}-P_f,&T_Q\dot Q_f&=Q_{inv}-Q_f.
\end{align}
The exact inverse of the dq power equations supplies the steady feedforward:
\begin{align}
i_{d,ff}&=\frac{v_dP_{ref}+v_qQ_{ref}}{D_V},&
i_{q,ff}&=\frac{v_qP_{ref}-v_dQ_{ref}}{D_V},\\
i_{d,r}^{raw}&=k_{pP}(P_{ref}-P_f)+k_{iP}\xi_P+i_{d,ff},\\
i_{q,r}^{raw}&=-k_{pQ}(Q_{ref}-Q_f)-k_{iQ}\xi_Q+i_{q,ff}.
\end{align}
The minus sign in the q-axis controller follows the generator convention
\(Q=v_qi_d-v_di_q\): at a locked PLL, positive reactive injection requires
negative \(i_q\).

\subsection{Current priority, balanced LVRT, and outer anti-windup}

In normal voltage, circular P-priority limits \((i_{d,r}^{raw},i_{q,r}^{raw})\)
to \(I_{max}\): \(i_d\) is clipped first and \(i_q\) receives the remaining
circle.  During the declared balanced positive-sequence voltage dip,
\begin{align}
i_{q,FRT}&=\clamp\!\left(K_{qv}\max(V_{dip}-|V|-d_b,0),-I_{qh1},I_{qh1}\right),\\
i_{d,LVPL}^{max}&=
\begin{cases}
0,&|V|\le Z_{erox},\\
L_{vpl1}\dfrac{|V|-Z_{erox}}{B_{rkpt}-Z_{erox}},&Z_{erox}<|V|<B_{rkpt},\\
L_{vpl1},&|V|\ge B_{rkpt},
\end{cases}\\
|i_d|&\le\min\!\left(i_{d,LVPL}^{max},\sqrt{I_{max}^2-i_q^2}\right).
\end{align}
This is reactive-current priority during the voltage dip; it does not cover
negative sequence or zero-voltage ride through.

After limiting, \(P_{lim}=v_di_{d,r}+v_qi_{q,r}\) and
\(Q_{lim}=v_qi_{d,r}-v_di_{q,r}\).  The outer integrators are
\begin{align}
\dot\xi_P&=AW_P(P_{ref}-P_f),&\dot\xi_Q&=AW_Q(Q_{ref}-Q_f).
\end{align}
Each \(AW\) returns zero only when its limiter is active and its error would
push the command farther outward; inward recovery remains active.  The
comparison tolerance is numerical metadata, not a tunable physical gain.

\subsection{Inner current PI, voltage command, and L-filter ODE}

With \(e_d=i_{d,r}-i_d\), \(e_q=i_{q,r}-i_q\), and
\(\omega_{PLL,pu}=1+\Delta\omega_{PLL}\),
\begin{align}
v_{td}^{raw}&=k_{pi}e_d+k_{ii}\xi_{id}+R_ti_d-\omega_{PLL,pu}Li_q+v_d,\\
v_{tq}^{raw}&=k_{pi}e_q+k_{ii}\xi_{iq}+R_ti_q+\omega_{PLL,pu}Li_d+v_q.
\end{align}
The vector \([v_{td}^{raw},v_{tq}^{raw}]\) is projected onto the modulation
circle of radius \(V_{t,max}=m_{max}V_{dc0}\).  The current PI integrators use
the same outward-hold/inward-release anti-windup rule:
\[
\dot\xi_{id}=AW_{id}(e_d),\qquad \dot\xi_{iq}=AW_{iq}(e_q).
\]
Because \(L\) is stored as per-unit reactance, the per-unit inductance in the
derivative denominator is \(L/\omega_b\):
\begin{align}
\dot i_d&=\frac{\omega_{PLL,pu}Li_q-R_ti_d+v_{td}-v_d}{L/\omega_b},\\
\dot i_q&=\frac{-\omega_{PLL,pu}Li_d-R_ti_q+v_{tq}-v_q}{L/\omega_b}.
\end{align}
Finally,
\[
I_{net}=\frac{(i_d+ji_q)e^{j\delta_{PLL}}}{\kappa},\qquad S=VI_{net}^*.
\]

\section{How PF, equilibrium, SSSA, and TS use the equations}

\subsection{Power flow}

PF does not integrate converter or SG states.  For each non-reference bus it
solves the specified active/reactive mismatch selected by the REF/PV/PQ
semantics.  In compact form, Newton iterations solve
\[
J_{PF}(z^{(m)})\Delta z^{(m)}=-F_{PF}(z^{(m)}),\qquad
z^{(m+1)}=z^{(m)}+\Delta z^{(m)},
\]
until \(\|F_{PF}\|_\infty\) is below tolerance.  The source-case REF/PV/PQ PF
is used only to freeze the constant-admittance load model.  The reported
operating point instead uses the mode-aware warm-start, changes every online
GFL terminal to its effective PQ semantics, and then accepts only the
subsequent coupled equilibrium satisfying the full device ODEs and every KCL
row.  Thus all reported bus, load, branch, and device tables share the same
voltage vector and SG1 plus four-GFL mode map.

\subsection{Mixed-device equilibrium}

Equilibrium uses the full device equations and solves
\[
\vect f_A(\vect x_0,\vect y_0,\vect u_0)=0,\qquad
\vect g(\vect x_0,\vect y_0,\vect u_0)=0,
\]
for active states and algebraic variables.  Frozen or inactive branch states
remain in the inventory but not in the Newton unknown vector.  GFL lock gives
\(v_q=0\), \(P_f=P_{ref}\), \(Q_f=Q_{ref}\), and stationary current states.
SG equilibrium solves constant \(T_m,E_{fd}\), which are then held in TS.

\subsection{Full-state SSSA}

Linearizing the same DAE at a fixed active regime gives
\[
\begin{bmatrix}\Delta\dot x\\0\end{bmatrix}=
\begin{bmatrix}f_x&f_y\\g_x&g_y\end{bmatrix}
\begin{bmatrix}\Delta x\\\Delta y\end{bmatrix},\qquad
A_{full}=f_x-f_y(g_y\backslash g_x).
\]
The published matrix is \(A=A_{full}(\mathcal A,\mathcal A)\), where
\(\mathcal A\) is the metadata-owned active-state map.  No eigenvalue is
manually deleted.  For \(Av_i=\lambda_i v_i\) and
\(A^Hu_i=\lambda_i^*u_i\), the vectors are paired and normalized so that
\(u_i^Hv_i=1\).  Signed participation is
\(p_{ki}=u_{ki}^*v_{ki}\); the displayed nonnegative ranking is
\(\rho_{ki}=|p_{ki}|/\sum_j|p_{ji}|\).  Mode number, raw eigen index, state
index, and device index are separate quantities.

\subsection{Time-domain simulation}

Between event transactions, the implicit trapezoidal step solves
\begin{align}
0&=x_{n+1}-x_n-\frac{h}{2}\left[f(x_n,y_n)+f(x_{n+1},y_{n+1})\right],\\
0&=g(x_{n+1},y_{n+1};\sigma_{n+1}).
\end{align}
The public sampling interval is \(h=0.01\) s.  This report uses an empty event
schedule, so the fixed all-GFL mode map remains unchanged for all 1500 steps.

\section{All-GFL equilibrium and SSSA results}

\begin{figure}[H]\centering
\includegraphics[width=\textwidth]{\datadir/matlab_normal_pf.pdf}
\caption{Accepted SG1 plus four-GFL equilibrium bus-voltage results.}
\end{figure}

\begin{figure}[H]\centering
\includegraphics[width=0.78\textwidth]{\datadir/matlab_normal_pf_convergence.pdf}
\caption{Mode-aware all-GFL PF warm-start convergence on a linear vertical scale.}
\end{figure}

\begin{table}[H]\centering\scriptsize
\caption{Solved bus voltages at the accepted SG1 plus four-GFL equilibrium.}
\pgfplotstabletypeset[col sep=comma,
 columns={bus_id,type,V_pu,V_kV,angle_deg},
 columns/bus_id/.style={column name=Bus},columns/type/.style={column name=Type,string type},
 columns/V_pu/.style={column name=$|V|$ pu,fixed,precision=4},
 columns/V_kV/.style={column name=$|V|$ kV,fixed,precision=2},
 columns/angle_deg/.style={column name=$\delta$ deg,fixed,precision=4},
 every head row/.style={before row=\toprule,after row=\midrule},every last row/.style={after row=\bottomrule}]
{\datadir/normal_pf_bus.csv}
\end{table}

\begin{table}[H]\centering\scriptsize
\caption{All-GFL equilibrium bus generation, voltage-dependent admittance
load, and net injection.  All quantities are on the 100-MVA system base and
positive net injection denotes power delivered to the network.}
\pgfplotstabletypeset[col sep=comma,
 columns={bus_id,P_gen_pu,Q_gen_pu,P_load_pu,Q_load_pu,P_net_pu,Q_net_pu},
 columns/bus_id/.style={column name=Bus},
 columns/P_gen_pu/.style={column name=$P_{\rm gen}$ pu,fixed,precision=4},
 columns/Q_gen_pu/.style={column name=$Q_{\rm gen}$ pu,fixed,precision=4},
 columns/P_load_pu/.style={column name=$P_{\rm load}$ pu,fixed,precision=4},
 columns/Q_load_pu/.style={column name=$Q_{\rm load}$ pu,fixed,precision=4},
 columns/P_net_pu/.style={column name=$P_{\rm net}$ pu,fixed,precision=4},
 columns/Q_net_pu/.style={column name=$Q_{\rm net}$ pu,fixed,precision=4},
 every head row/.style={before row=\toprule,after row=\midrule},every last row/.style={after row=\bottomrule}]
{\datadir/normal_pf_power.csv}
\end{table}

\begin{table}[H]\centering\scriptsize
\caption{All-GFL equilibrium branch sending-end flows and losses.  Power
quantities are on the 100-MVA system base; the sign follows the listed
from--to orientation.}
\pgfplotstabletypeset[col sep=comma,
 columns={line,from_bus,to_bus,P_from_pu,Q_from_pu,P_loss_pu,Q_loss_pu},
 columns/line/.style={column name=Line},columns/from_bus/.style={column name=From},
 columns/to_bus/.style={column name=To},
 columns/P_from_pu/.style={column name=$P_{ij}$ pu,fixed,precision=4},
 columns/Q_from_pu/.style={column name=$Q_{ij}$ pu,fixed,precision=4},
 columns/P_loss_pu/.style={column name=$P_{\rm loss}$ pu,fixed,precision=6},
 columns/Q_loss_pu/.style={column name=$Q_{\rm loss}$ pu,fixed,precision=6},
 every head row/.style={before row=\toprule,after row=\midrule},every last row/.style={after row=\bottomrule}]
{\datadir/normal_pf_line_flow.csv}
\end{table}

\begin{table}[H]\centering\scriptsize
\caption{Solved device injections at the same all-GFL mixed-device equilibrium.}
\pgfplotstabletypeset[col sep=comma,
 columns={device_index,device_id,bus_id,mode,P_MW,Q_MVAr,V_pu,V_kV},
 columns/device_index/.style={column name=Index},columns/device_id/.style={column name=Device,string type},
 columns/bus_id/.style={column name=Bus},columns/mode/.style={column name=Mode,string type},
 columns/P_MW/.style={column name=$P$ MW,fixed,precision=3},
 columns/Q_MVAr/.style={column name=$Q$ MVAr,fixed,precision=3},
 columns/V_pu/.style={column name=$|V|$ pu,fixed,precision=4},
 columns/V_kV/.style={column name=$|V|$ kV,fixed,precision=3},
 every head row/.style={before row=\toprule,after row=\midrule},every last row/.style={after row=\bottomrule}]
{\datadir/normal_resource_pf.csv}
\end{table}

\begin{figure}[H]\centering
\includegraphics[width=\textwidth]{\datadir/matlab_normal_sssa.pdf}
\caption{All 45 SG1 plus four-GFL eigenvalues (left) and the low/medium-frequency
detail (right).  Each point is one eigenvalue of the same 45-state reduced
matrix.  The dashed line is \(\Re\lambda=0\); roots to its right are unstable.}
\end{figure}

\subsection{SG1 plus four GFL--RMS10 normal-operation SSSA}

The explicit all-GFL mode map keeps SG1 online and assigns IBR2, IBR3, IBR6,
and IBR8 to GFL--RMS10.  Its independently solved equilibrium has 45 active
states.  The complete spectrum below is published without filtering: the four
large positive real PLL roots and the small positive SG angle root make the
physical result \textbf{UNSTABLE}.  Successful equilibrium and SSSA execution
therefore must not be interpreted as a stable classification.

\begin{table}[p]\centering\tiny
\caption{Complete 45-root SSSA result for SG1 plus four GFL--RMS10 resources.}
\pgfplotstabletypeset[col sep=comma,
 columns={mode,raw_eigen_index,real_1_per_s,imag_1_per_s,frequency_Hz,damping_ratio,dominant_device,local_state,participation_percent},
 columns/mode/.style={column name=No.},columns/raw_eigen_index/.style={column name=Raw},
 columns/real_1_per_s/.style={column name=$\Re\lambda$ (s$^{-1}$),sci,sci zerofill,precision=3},
 columns/imag_1_per_s/.style={column name=$\Im\lambda$ (s$^{-1}$),sci,sci zerofill,precision=3},
 columns/frequency_Hz/.style={column name=$f$ (Hz),sci,sci zerofill,precision=3},
 columns/damping_ratio/.style={column name=$\zeta$,fixed,precision=4},
 columns/dominant_device/.style={column name=Device,string type},
 columns/local_state/.style={column name=Local state},
 columns/participation_percent/.style={column name=Part. (\%),fixed,precision=2},
 every head row/.style={before row=\toprule,after row=\midrule},every last row/.style={after row=\bottomrule}]
{\datadir/normal_sssa_eigenvalues.csv}
\end{table}

\clearpage
\section{Event-free time-domain result}

The SG1 plus four-GFL normal-operation TS starts from the same converged
equilibrium used by SSSA and completes all 1500 fixed 0.01-s steps to 15 s.
The empty event schedule keeps SG1 online and all four IBRs in GFL-RMS10 mode.
No fault, breaker transaction, or automatic mode switch is included.

\begin{figure}[H]\centering
\includegraphics[width=\textwidth]{\datadir/matlab_normal_event_free_ts.pdf}
\caption{MATLAB event-free TS for SG1 plus four GFL-RMS10 resources over 15 s.}
\end{figure}

\section{Comparison with the conventional SG analysis path}

The conventional multi-SG path and the mixed IBR path share the network and
analysis mathematics, but not the device states.  A classical SG contributes
\([\delta,\omega]\), while the present SG contributes the EMF6 inventory and
the converter models contribute control and physical current states.  PF is
common because all resources appear as scheduled injections and voltage
constraints.  Equilibrium, SSSA, and TS are generic consumers of the device
ABI: they do not contain a device-specific eigenvalue matrix.  In the present
fixed mode map, the four IBRs contribute only their GFL-RMS10 coordinates and
current laws to the Schur-complement and implicit-trapezoidal algorithms.

\section{Limitations and interpretation}

\begin{itemize}
\item The GFL fault contract is balanced positive-sequence RMS LVRT.  It is
not an EMT, negative-sequence, or zero-voltage model.
\item A completed numerical SSSA execution is distinct from a stable physical
classification.  Positive-real-part eigenvalues remain in the table and plot.
\item The documented grid-forming branch remains inactive in every numerical
result in this report; its equations are model documentation, not evidence that
any IBR was assigned that mode.
\end{itemize}

\section{Conclusion}

The report exposes a traceable chain from case data through device equations,
state order, initialization, DAE assembly, PF, SSSA, and TS.  The same
nonlinear SG and four GFL-RMS10 equations initialize the equilibrium, are
differentiated for SSSA, and are integrated for TS.  Native \LaTeX{} plots
retain the exact MATLAB series while enforcing common typography and symmetric
layout.  The inactive grid-forming equations remain documented without being
mixed into the reported numerical operating point.

\begin{thebibliography}{9}
\bibitem{kodsi} S. K. M. Kodsi and C. A. Ca\~nizares,
``Modeling and Simulation of IEEE 14 Bus System with FACTS Controllers,''
University of Waterloo Technical Report 2003-3, Table A.2, 2003.
\bibitem{ieee1110} IEEE Std 1110-2002, \emph{IEEE Guide for Synchronous
Generator Modeling Practices and Applications in Power System Stability Analyses}.
\bibitem{nrel90260} NREL, \emph{REGFM\_B1: An Electromagnetic Transient Model
for Grid-Forming Inverters}, NREL/TP-5D00-90260, 2024.
\bibitem{yazdani} A. Yazdani and R. Iravani, \emph{Voltage-Sourced Converters
in Power Systems}, Wiley, 2010.
\bibitem{teodorescu} R. Teodorescu, M. Liserre, and P. Rodr\'iguez,
\emph{Grid Converters for Photovoltaic and Wind Power Systems}, Wiley, 2011.
\bibitem{bacha} S. Bacha, I. Munteanu, and A. I. Bratcu,
\emph{Power Electronic Converters Modeling and Control}, Springer, 2014,
doi:10.1007/978-1-4471-5478-5.
\end{thebibliography}

\end{document}
